\vspace{2.5pt plus 3pt minus 0.5pt}\noindent\hspace*{1.2em}\kbvnum{(9)}\hspace{0.55em}ಸರ್ವನಾಮಗಳು ಎಂದರೇನು ? ಉದಾಹರಿಸಿ.\par

\vspace{2.5pt plus 3pt minus 0.5pt}\noindent\hspace*{1.2em}\kbvnum{(10)}\hspace{0.55em}ರಾಮನು ದಶರಥನ ಮಗ. ರಾಮನು ತಂದೆಯ ಮಾತಿನಂತೆ ಕಾಡಿಗೆ ಹೋದನು. ರಾಮನೊಂದಿಗೆ ಸೀತೆಯೂ ಕಾಡಿಗೆ ಹೋದಳು. ಸೀತಾರಾಮರ ಜೊತೆಯಲ್ಲಿ ಲಕ್ಷ್ಮಣನೂ ಹೊರಟನು.\par

\indent — ಈ ವಾಕ್ಯಗಳಲ್ಲಿ ಅಗತ್ಯವಿದ್ದೆಡೆ ಸರ್ವನಾಮಗಳನ್ನು ಸೇರಿಸಿ ಬರೆಯಿರಿ.\par

\vspace{2.5pt plus 3pt minus 0.5pt}\noindent\hspace*{1.2em}\kbvnum{(11)}\hspace{0.55em}ಉತ್ತಮ ಪುರುಷ ಸರ್ವನಾಮಗಳಿಗೂ ಮಧ್ಯಮ ಪುರುಷ ಸರ್ವನಾಮಗಳಿಗೂ ಇರುವ ವ್ಯತ್ಯಾಸವೇನು ? ಉದಾಹರಿಸಿ.\par

\vspace{2.5pt plus 3pt minus 0.5pt}\noindent\hspace*{1.2em}\kbvnum{(12)}\hspace{0.55em}ನಿನ್ನ ವಾಸಸ್ಥಳ ಯಾವುದು? ನೀನು ಎಲ್ಲಿಂದ ಬಂದೆ? ನಿನ್ನನ್ನು ಕಳುಹಿಸಿದವರು ಯಾರು? ನೀನು ಬಂದ ಉದ್ದೇಶವೇನು?\par

\indent — ಇಲ್ಲಿ ಕಾಣಿಸಿರುವ ಪ್ರಶ್ನಾರ್ಥಕ ಸರ್ವನಾಮಗಳನ್ನು ಗುರ್ತಿಸಿ.\par

\kbvchapter{ಪಾಠ ೯}{ನಾಮಪದಗಳ ರೂಪ ಬದಲಾವಣೆ ({\kbvlatin i})}

\kbvsubtitle{ಲಿಂಗ}

\noindent\hspace*{2.2em}ಭೀಮಣ್ಣ ಓಡುತ್ತಿದ್ದಾನೆ.\par

\noindent\hspace*{2.2em}ಆ ಕಮಲೆ ಹಾಡುತ್ತಿದ್ದಾಳೆ.\par

\noindent\hspace*{2.2em}ನವಿಲು ಕುಣಿಯುತ್ತಿದೆ.\par

\noindent\hspace*{2.2em}ಗಡಿಯಾರ ನಿಂತುಹೋಗಿದೆ.\par

\indent ಈ ವಾಕ್ಯಗಳಲ್ಲಿ ದಪ್ಪ ಅಕ್ಷರಗಳಲ್ಲಿರುವ ಪದಗಳನ್ನು ಗಮನಿಸಿ. “ಭೀಮಣ್ಣ” ಎಂಬ ಶಬ್ದ ಕಿವಿಗೆ ಬಿದ್ದೊಡನೆ ಇವನು ಪುರುಷ {[}ಗಂಡಸು{]} ಎಂದು ಹೇಳಬಹುದು. ಈ ಕಾರಣದಿಂದ “ಭೀಮಣ್ಣ” ಎನ್ನುವುದು ಪುರುಷಲಿಂಗ ಶಬ್ದ ಎನ್ನುತ್ತೇವೆ. “ಕಮಲೆ” - ಎಂಬ ಶಬ್ದ ಕಿವಿಗೆ ಬಿದ್ದೊಡನೆ ಇವಳು ಸ್ತ್ರೀ {[}ಹೆಂಗಸು{]} ಎಂದು ಹೇಳಬಹುದು. ಈ ಕಾರಣದಿಂದ “ಕಮಲೆ” ಎನ್ನುವುದು ಸ್ತ್ರೀಲಿಂಗ ಶಬ್ದ ಎನ್ನುತ್ತೇವೆ. ನವಿಲು, ಗಡಿಯಾರ - ಈ ಶಬ್ದಗಳು ಮನಸ್ಸಿನ ಮೇಲೆ ಗಂಡು, ಎಂಬ ಅಭಿಪ್ರಾಯವನ್ನಾಗಲೀ ಹೆಣ್ಣು ಎಂಬ ಅಭಿಪ್ರಾಯವನ್ನಾಗಲೀ ಉಂಟುಮಾಡುವುದಿಲ್ಲ. ಆದುದರಿಂದ “ನವಿಲು”, “ಗಡಿಯಾರ” ಎಂಬ ಶಬ್ದಗಳನ್ನು ನಪುಂಸಕ ಲಿಂಗ ಶಬ್ದಗಳೆನ್ನುತ್ತೇವೆ.\par
